\documentclass[11pt,a4paper]{article}

\usepackage[T1]{fontenc}
\usepackage[utf8]{inputenc}
\usepackage{libertinus}
\usepackage{microtype}
\usepackage[a4paper,margin=27mm,headheight=14pt]{geometry}
\usepackage{mathtools,amssymb,amsthm,bm}
\usepackage{booktabs,longtable,array,tabularx}
\usepackage{enumitem}
\usepackage{xcolor}
\usepackage[most]{tcolorbox}
\usepackage{listings}
\usepackage{fancyhdr}
\usepackage{titlesec}
\usepackage{hyperref}
\usepackage[nameinlink,capitalize,noabbrev]{cleveref}

\definecolor{inkblue}{RGB}{25,55,92}
\definecolor{softblue}{RGB}{238,244,250}
\definecolor{softgray}{RGB}{246,247,249}
\definecolor{rulegray}{RGB}{110,118,128}
\hypersetup{
  colorlinks=true,
  linkcolor=inkblue,
  citecolor=inkblue,
  urlcolor=inkblue,
  pdftitle={The Full Saddle--Transseries Expansion of the Bell Numbers},
  pdfauthor={Prepared for the ProveIt research archive},
  pdfsubject={Bell-number asymptotics, saddle-point coefficient calculus, Lambert W transseries}
}

\pagestyle{fancy}
\fancyhf{}
\fancyhead[L]{\small\textsc{Bell-number transseries}}
\fancyhead[R]{\small\nouppercase{\leftmark}}
\fancyfoot[C]{\thepage}
\renewcommand{\headrulewidth}{0.4pt}
\titleformat{\section}{\Large\bfseries\color{inkblue}}{\thesection}{0.75em}{}
\titleformat{\subsection}{\large\bfseries\color{inkblue}}{\thesubsection}{0.7em}{}
\titleformat{\subsubsection}{\normalsize\bfseries}{\thesubsubsection}{0.6em}{}

\newtheorem{theorem}{Theorem}[section]
\newtheorem{proposition}[theorem]{Proposition}
\newtheorem{lemma}[theorem]{Lemma}
\newtheorem{corollary}[theorem]{Corollary}
\newtheorem{definition}[theorem]{Definition}
\newtheorem{remark}[theorem]{Remark}
\newtheorem{example}[theorem]{Example}
\theoremstyle{definition}
\newtheorem{algorithm}[theorem]{Algorithm}

\newtcolorbox{scopebox}{
  colback=softblue,
  colframe=inkblue,
  boxrule=0.7pt,
  arc=1.5mm,
  left=3mm,right=3mm,top=2.5mm,bottom=2.5mm,
  title=Scope and truth status,
  fonttitle=\bfseries
}
\newtcolorbox{resultbox}[1][]{
  colback=softgray,
  colframe=rulegray,
  boxrule=0.5pt,
  arc=1mm,
  left=3mm,right=3mm,top=2.5mm,bottom=2.5mm,
  #1
}

\lstdefinestyle{code}{
  language=Python,
  basicstyle=\ttfamily\footnotesize,
  keywordstyle=\bfseries,
  commentstyle=\itshape,
  showstringspaces=false,
  breaklines=true,
  frame=single,
  rulecolor=\color{rulegray},
  backgroundcolor=\color{softgray},
  columns=fullflexible,
  keepspaces=true,
  aboveskip=0.8em,
  belowskip=0.8em
}

\newcommand{\Bell}{\mathrm B}
\newcommand{\Touch}{\mathrm T}
\newcommand{\W}{\mathrm W}
\newcommand{\ee}{\mathrm e}
\newcommand{\ii}{\mathrm i}
\newcommand{\dd}{\,\mathrm d}
\newcommand{\Gaus}{\mathfrak G}
\newcommand{\cB}{\mathcal B}
\newcommand{\cA}{\mathcal A}
\newcommand{\cE}{\mathcal E}
\newcommand{\cJ}{\mathcal J}
\newcommand{\cL}{\mathcal L}
\newcommand{\Log}{\operatorname{Log}}
\newcommand{\Res}{\operatorname{Res}}
\newcommand{\bigO}{\mathcal O}
\newcommand{\smallo}{o}
\newcommand{\ind}{\mathbf 1}
\newcommand{\Stwo}[2]{\genfrac{\{}{\}}{0pt}{}{#1}{#2}}
\newcommand{\Sone}[2]{\genfrac{[}{]}{0pt}{}{#1}{#2}}
\DeclareMathOperator{\RePart}{Re}
\DeclareMathOperator{\ImPart}{Im}
\DeclareMathOperator{\length}{len}
\DeclarePairedDelimiter{\abs}{\lvert}{\rvert}
\DeclarePairedDelimiter{\norm}{\lVert}{\rVert}
\DeclarePairedDelimiter{\paren}{(}{)}
\DeclarePairedDelimiter{\bracks}{[}{]}
\DeclarePairedDelimiter{\set}{\{} {\}}

\setlist[itemize]{leftmargin=1.6em,itemsep=0.25em,topsep=0.4em}
\setlist[enumerate]{leftmargin=1.8em,itemsep=0.3em,topsep=0.4em}
\allowdisplaybreaks[2]

\title{\vspace{-1.0cm}\textbf{The Full Saddle--Transseries Expansion\\of the Bell Numbers}\\[0.45em]
\large Lambert--$W$ sectors and combinatorial formulae for every coefficient}
\author{Prepared for the ProveIt research archive}
\date{September 2026}

\begin{document}
\maketitle

\begin{abstract}
Let $\Bell_n$ be the $n$th Bell number and let $r=\W_0(n)$, so that
$r\ee^r=n$.  We derive a complete fixed-order asymptotic expansion about the
principal Cauchy saddle,
\[
 \Bell_n=
 \frac{n!\exp(\ee^r-1)}{r^n\sqrt{2\pi n(1+r)}}
 \left\{\sum_{m=0}^{M-1}C_m(r)\left(\frac rn\right)^m
 +\bigO_M\!\left(\left(\frac rn\right)^M\right)\right\},
\]
and give four equivalent, finite descriptions of every coefficient $C_m(r)$:
a Gaussian coefficient extractor, an exponential-Bell-polynomial formula, a
triangular recurrence, and an explicit sum over the integer partitions of
$2m$.  In particular,
$C_m(r)=P_m(r)/\{r^m(1+r)^{3m}\}$ with
$P_m\in\mathbb Q[r]$ and $\deg P_m\le 4m$.  The first five coefficients are
written out, and the remainder is proved by a central-arc expansion together
with a uniform complementary-arc estimate.

We then interface this expansion with Stirling's series, derive an all-order
factorial-free form and an all-order logarithmic (cumulant) form, and identify
$\lim_{r\to\infty}C_m(r)$ with the reciprocal Stirling coefficients.  The
Lambert function is itself expanded on every branch by explicit Stirling-first-
kind polynomials, converting the answer into the canonical nested
$\log n$--$\log\log n$ transseries.

Finally, using the Hankel representation with $N=n+1$, we retain every
contributory steepest-descent saddle $w_k=\W_k(N)$.  Each thimble sector is
expanded to arbitrary order, with a second finite partition formula for its
coefficients.  The original Bell contour fixes the Stokes multipliers, while
the singulants show that the $k\ne0$ sectors are exponentially small like
$\exp\{-2\pi^2k^2N/\W_0(N)^2\}$ for fixed $k$.  Thus the result includes both
the complete perturbative series and the beyond-all-orders Lambert--$W$
sectors, with explicit algorithms and numerical checks.
\end{abstract}

\begin{scopebox}
The word \emph{full} is used in the following precise sense.  First, every
coefficient in the principal Poincar\'e expansion is given by a finite formula.
Second, every complex Lambert--$W$ saddle has an all-order local expansion with
a finite coefficient formula.  Third, the Bell contour selects a finite set of
active thimbles at each fixed value of the large parameter, and this selection
is stated explicitly in terms of Stokes data.  Fourth, every Lambert branch can
be re-expanded into the polynomial--logarithmic transseries scale.

The principal fixed-order error estimate in this article is analytic and
uniform for sufficiently large $n$ (with the truncation order fixed).  The
multisaddle expansion is sectorwise: no claim of a single error bound uniform
in a branch index growing with $n$ is made.  At a Stokes threshold, lateral or
median interpretation is required.  The derivations are human-readable and
computer-algebra checked, but are not claimed to be machine-verified in a proof
assistant.
\end{scopebox}

\tableofcontents
\bigskip

\section{Introduction}

The Bell number $\Bell_n$ counts the set partitions of an $n$-element labelled
set.  Its exponential generating function is
\begin{equation}\label{eq:egf}
 \exp(\ee^z-1)=\sum_{n\ge0}\Bell_n\frac{z^n}{n!}.
\end{equation}
The coefficient problem is classical, but its complete asymptotic structure is
richer than the familiar leading formula suggests.  The positive saddle moves
to infinity with $n$ and is determined implicitly by
\[
 r\ee^r=n,
 \qquad r=\W_0(n).
\]
Consequently the natural small parameter is not $1/n$ alone, but
\[
 q:=\ee^{-r}=\frac rn\sim\frac{\log n}{n}.
\]
At the same time a Hankel representation exposes infinitely many complex
saddles $\W_k(n+1)$; a finite, contour-dependent subset contributes for each
fixed $n$.  The nonprincipal contributions lie beyond every power of $q$ and
therefore belong to a genuine transseries rather than to the principal
Poincar\'e series.

The classical analyses of Moser--Wyman, Hayman, de Bruijn, and later authors
establish the leading shape and various refinements.  Paris's steepest-descent
analysis of the Touchard polynomials makes the Lambert-branch saddle geometry
and its Stokes changes explicit.  Hwang identified a striking limiting
connection between saddle coefficients and Stirling's formula.  The purpose of
the present article is to assemble these ingredients into a self-contained
coefficient calculus that satisfies five obligations simultaneously:
\begin{enumerate}
 \item identify the actual asymptotic scale;
 \item give a finite, executable formula for every coefficient;
 \item state the exact shift convention ($n$ versus $n+1$);
 \item prove the central-arc remainder estimate and a complementary-arc bound;
 \item include the exponentially small Lambert-branch sectors and their Stokes
       selection.
\end{enumerate}
These are also the obligations singled out in the audit notes of the
\emph{Combinatorial Coefficient Calculus} draft in the ProveIt repository.  The
formal/analytic separation and the branch-aware logarithmic expansions follow
the conventions of the companion \emph{series-and-transseries} drafts.

\subsection{Summary of the principal result}

For orientation, define the leading saddle approximation
\begin{equation}\label{eq:leading-def}
 L_n(r):=
 \frac{n!\exp(\ee^r-1)}{r^n\sqrt{2\pi n(1+r)}},
 \qquad r=\W_0(n).
\end{equation}
Then for every fixed $M\ge1$,
\begin{equation}\label{eq:intro-main}
 \Bell_n=L_n(r)
 \left(\sum_{m=0}^{M-1}C_m(r)q^m+\bigO_M(q^M)\right),
 \qquad q=\frac rn.
\end{equation}
The coefficient $C_m$ is a rational function with a tightly controlled
denominator.  If $\lambda=(1^{m_1}2^{m_2}\cdots)$ is a partition of $2m$,
write $\ell(\lambda)=\sum_jm_j$.  Let
\begin{equation}\label{eq:touchard-def-intro}
 \Touch_j(x):=\sum_{a=0}^j \Stwo{j}{a}x^a
\end{equation}
be the $j$th Touchard polynomial.  The finite coefficient formula is
\begin{equation}\label{eq:intro-partition}
 \boxed{
 C_m(r)=\sum_{\lambda\vdash 2m}
 \frac{(-1)^{m+\ell(\lambda)}
       (2m+2\ell(\lambda)-1)!!}
      {[r(1+r)]^{m+\ell(\lambda)}}
 \prod_{j\ge1}
 \frac{\Touch_{j+2}(r)^{m_j}}
      {m_j!\,((j+2)!)^{m_j}}.}
\end{equation}
For $m=0$ the empty partition gives $C_0=1$.  Formula
\eqref{eq:intro-partition} is the central combinatorial object of the article.

\subsection{What is new in the organization}

No claim is made that every ingredient below is historically new.  The
contribution is the explicit synthesis and normalization:
\begin{itemize}
 \item the circle-saddle expansion is placed directly in the scale
       $q=r/n$, with a finite partition formula and a fixed-order remainder;
 \item the denominator and degree structure of all $C_m$ is proved;
 \item the Stirling, logarithmic, and polynomial--logarithmic interfaces are
       made algorithmic;
 \item the same partition calculus is transported to all Hankel thimbles,
       yielding an explicit full saddle transseries;
 \item the relation between the $r=\W_0(n)$ and $w=\W_0(n+1)$ normalizations is
       kept visible, avoiding a common off-by-one ambiguity.
\end{itemize}

\section{Exact representations and notation}

\subsection{Bell and Touchard polynomials}

The Touchard polynomials are
\begin{equation}\label{eq:touchard}
 \Touch_n(x)=\ee^{-x}\sum_{j\ge0}\frac{j^n x^j}{j!}
 =\sum_{a=0}^n \Stwo{n}{a}x^a,
\end{equation}
so $\Bell_n=\Touch_n(1)$.  Their exponential generating function is
\begin{equation}\label{eq:touchard-egf}
 \exp\!\bigl(x(\ee^z-1)\bigr)
 =\sum_{n\ge0}\Touch_n(x)\frac{z^n}{n!}.
\end{equation}
The same polynomials occur as logarithmic derivatives of the exponential:
\begin{equation}\label{eq:derivative-touchard}
 (xD_x)^j\ee^x=\ee^x\Touch_j(x),
 \qquad D_x:=\frac{\dd}{\dd x}.
\end{equation}
This identity is exactly what converts the local saddle derivatives into
Stirling-number data.

The first few polynomials are
\begin{align*}
 \Touch_0(x)&=1, &
 \Touch_1(x)&=x,\\
 \Touch_2(x)&=x^2+x, &
 \Touch_3(x)&=x^3+3x^2+x,\\
 \Touch_4(x)&=x^4+6x^3+7x^2+x, &
 \Touch_5(x)&=x^5+10x^4+25x^3+15x^2+x.
\end{align*}

\subsection{The circle representation}

Cauchy's coefficient formula applied to \eqref{eq:egf}, on the circle
$z=r\ee^{\ii\theta}$, gives the exact identity
\begin{equation}\label{eq:circle}
 \frac{\Bell_n}{n!}
 =\frac{\exp(\ee^r-1)}{2\pi r^n}
 \int_{-\pi}^{\pi}\exp\!\bigl(\Phi_r(\theta)\bigr)\dd\theta,
\end{equation}
where
\begin{equation}\label{eq:phase-circle}
 \Phi_r(\theta):=\ee^{r\ee^{\ii\theta}}-\ee^r-\ii n\theta.
\end{equation}
The saddle equation $\Phi_r'(0)=0$ is $r\ee^r=n$.  Throughout the principal
analysis we therefore set
\begin{equation}\label{eq:rqa}
 r:=\W_0(n),\qquad q:=\ee^{-r}=\frac rn,
 \qquad A:=n(1+r)=\frac{r(1+r)}q.
\end{equation}
Here $A=-\Phi_r''(0)$ is the Gaussian curvature.

\subsection{The Hankel representation and the shift convention}

A second representation is needed to see all complex saddles.  Put
\begin{equation}\label{eq:N-shift}
 N:=n+1.
\end{equation}
By inserting a Hankel formula for the powers in Dobinski's sum, or directly
from \eqref{eq:touchard-egf}, one obtains
\begin{equation}\label{eq:hankel}
 \Bell_n=\Touch_{N-1}(1)
 =\frac{\Gamma(N)\ee^{-1}}{2\pi\ii}
 \int_{-\infty}^{(0+)}\exp\!\bigl(N\psi_N(t)\bigr)\dd t,
 \qquad
 \psi_N(t):=\frac{\ee^t}{N}-\Log t.
\end{equation}
The loop starts at $-\infty$, encircles the origin counterclockwise, and
returns to $-\infty$; the logarithm is cut along the negative real axis.  Its
saddles satisfy
\begin{equation}\label{eq:hankel-saddles}
 t\ee^t=N,
 \qquad w_k:=\W_k(N),\quad k\in\mathbb Z.
\end{equation}
The shift in \eqref{eq:N-shift} is essential: the circle normalization uses
$\W_0(n)$, while the Touchard/Hankel normalization for $\Bell_n=\Touch_n(1)$
uses $\W_k(n+1)$.

\subsection{Asymptotic notation}

All $\bigO_M$ constants may depend on the fixed truncation order $M$ but not on
$n$.  Unless a branch index is explicitly allowed to vary, a statement about
$w_k$ is for fixed $k$ as $N\to\infty$.  Square roots in complex saddle
formulae are fixed by continuous transport along the oriented
steepest-descent thimble, with the conjugacy convention
$\cJ_{-k}=\overline{\cJ_k}$ for $N>0$.

\section{The principal circle saddle to all orders}

\subsection{Exact local phase}

Differentiating \eqref{eq:phase-circle} and using
\eqref{eq:derivative-touchard} gives
\begin{equation}\label{eq:phase-taylor}
 \Phi_r(\theta)
 =-\frac A2\theta^2+
 \sum_{j\ge3}\frac{\ee^r\Touch_j(r)}{j!}(\ii\theta)^j.
\end{equation}
Scale
\begin{equation}\label{eq:scale-u}
 \theta=\frac{u}{\sqrt A}
 =\sqrt{\frac{q}{r(1+r)}}\,u.
\end{equation}
For $j\ge3$ define
\begin{equation}\label{eq:alpha}
 \alpha_j(r):=
 \frac{\Touch_j(r)}{j!\,[r(1+r)]^{j/2}}.
\end{equation}
Then
\begin{equation}\label{eq:scaled-phase}
 \Phi_r\!\left(\frac{u}{\sqrt A}\right)
 =-\frac{u^2}{2}
 +\sum_{j\ge3}\alpha_j(r)(\ii u)^j q^{(j-2)/2}.
\end{equation}
Since $\Touch_j(r)$ has nonnegative coefficients and degree $j$,
\begin{equation}\label{eq:alpha-bound}
 0\le \alpha_j(r)\le\frac{\Bell_j}{j!}
 \qquad (r\ge1),
\end{equation}
so every fixed family of normalized derivatives is uniformly bounded.

\subsection{Main fixed-order theorem}

\begin{theorem}[Complete principal expansion]\label{thm:principal}
Let $M\ge1$ be fixed, $r=\W_0(n)$, and $q=r/n$.  As $n\to\infty$,
\begin{equation}\label{eq:principal-theorem}
 \Bell_n=
 \frac{n!\exp(\ee^r-1)}{r^n\sqrt{2\pi n(1+r)}}
 \left\{\sum_{m=0}^{M-1}C_m(r)q^m+R_M(n)\right\},
\end{equation}
where
\begin{equation}\label{eq:principal-rem}
 R_M(n)=\bigO_M(q^M).
\end{equation}
More precisely, the part of the contour outside a fixed small neighbourhood
of $\theta=0$ is
\begin{equation}\label{eq:outer-exp-small}
 \bigO\!\left(\exp\!\left\{-c\frac{n}{r^2}\right\}\right)
\end{equation}
relative to the unscaled central integral, for some absolute $c>0$.  The
coefficients $C_m(r)$ are given in \cref{thm:coefficient-calculus} below.
\end{theorem}

\begin{proof}
We give the proof in enough detail to identify the uniformity.  Write
$t=q^{1/2}$ and choose a fixed sufficiently small $\eta>0$.  Split the circle
into three regions:
\begin{align*}
 \mathcal C_0&:\ |u|\le t^{-1/4},\\
 \mathcal C_1&:\ t^{-1/4}<|u|\le \eta\sqrt A/r,\\
 \mathcal C_2&:\ \eta/r<|\theta|\le\pi,
\end{align*}
where $u=\sqrt A\,\theta$.  Because $\sqrt A/r\to\infty$, these regions are
well ordered for large $n$.

On $\mathcal C_0$, \eqref{eq:scaled-phase} and
\eqref{eq:alpha-bound} imply that the nonquadratic exponent is
$\bigO(t|u|^3)=\bigO(t^{1/4})$.  Expanding its exponential through degree
$2M-1$ in $t$ gives
\begin{equation}\label{eq:central-expansion-proof}
 \exp\!\left(\sum_{\ell\ge1}h_\ell(u;r)t^\ell\right)
 =\sum_{\nu=0}^{2M-1}P_\nu(u;r)t^\nu
 +\bigO_M\!\left(t^{2M}(1+|u|)^{6M}\ee^{u^2/4}\right),
\end{equation}
with $h_\ell$ and $P_\nu$ defined in
\eqref{eq:h-def}--\eqref{eq:P-generating}.  After multiplication by
$\ee^{-u^2/2}$, the error is integrable and contributes $\bigO_M(t^{2M})$.
The interval is symmetric, $P_\nu$ has parity $\nu$, and hence every odd
coefficient integrates to zero.  Extending the resulting Gaussian integrals
to $\mathbb R$ changes them by
$\bigO(\exp\{-c t^{-1/2}\})$.

For $|\theta|\le\eta/r$, Taylor's theorem applied to
$\RePart\Phi_r(\theta)$, with $\eta$ chosen small, yields the uniform quadratic
bound
\begin{equation}\label{eq:local-decay}
 \RePart\Phi_r(\theta)\le-\frac14A\theta^2.
\end{equation}
Thus $\mathcal C_1$ contributes
$\bigO(\exp\{-c t^{-1/2}\})$.

Finally,
\begin{equation}\label{eq:global-real-bound}
 \RePart\Phi_r(\theta)
 =\ee^{r\cos\theta}\cos(r\sin\theta)-\ee^r
 \le \ee^{r\cos\theta}-\ee^r.
\end{equation}
For $\eta/r\le|\theta|\le\pi$, the inequality
$1-\cos\theta\ge2\theta^2/\pi^2$ gives
\begin{align*}
 \RePart\Phi_r(\theta)
 &\le-\ee^r\left(1-\exp\{-r(1-\cos\theta)\}\right)\\
 &\le-c\frac{\ee^r}{r}
 =-c\frac{n}{r^2},
\end{align*}
which proves \eqref{eq:outer-exp-small}.  Reinstating the Jacobian
$\dd\theta=\dd u/\sqrt A$ and the factor in \eqref{eq:circle} gives
\eqref{eq:principal-theorem}.  Since $t^{2M}=q^M$ and all discarded contour
pieces are smaller than every power of $q$, \eqref{eq:principal-rem} follows.
\end{proof}

\begin{remark}[Genuine scale]\label{rem:genuine-scale}
The expansion parameter is $q=r/n$, not $1/n$ with bounded coefficients.  The
coefficients $C_m(r)$ remain bounded as $r\to\infty$, whereas writing the same
series directly in powers of $1/n$ produces coefficients of order $r^m$.
This distinction is indispensable when stating a correct remainder.
\end{remark}

\section{Combinatorial coefficient calculus}

\subsection{The Gaussian functional}

Define the normalized Gaussian functional
\begin{equation}\label{eq:gaussian-functional}
 \Gaus[p]:=\frac1{\sqrt{2\pi}}
 \int_{-\infty}^{\infty}\ee^{-u^2/2}p(u)\dd u.
\end{equation}
Thus
\begin{equation}\label{eq:gaussian-moments}
 \Gaus[u^{2a}]=(2a-1)!!,
 \qquad \Gaus[u^{2a+1}]=0.
\end{equation}
For $\ell\ge1$, put
\begin{equation}\label{eq:h-def}
 h_\ell(u;r):=
 \alpha_{\ell+2}(r)(\ii u)^{\ell+2}.
\end{equation}
Define polynomials $P_\nu(u;r)$ by
\begin{equation}\label{eq:P-generating}
 \exp\!\left(\sum_{\ell\ge1}h_\ell(u;r)t^\ell\right)
 =\sum_{\nu\ge0}P_\nu(u;r)t^\nu.
\end{equation}

\begin{theorem}[Equivalent formulae for every coefficient]
\label{thm:coefficient-calculus}
For $m\ge0$, the coefficient in \cref{thm:principal} has each of the following
finite descriptions.

\smallskip
\noindent\textup{(i) Gaussian coefficient extraction:}
\begin{equation}\label{eq:C-gaussian}
 C_m(r)=\Gaus[P_{2m}(u;r)]
 =\Gaus\!\left([t^{2m}]
 \exp\!\left\{\sum_{\ell\ge1}h_\ell(u;r)t^\ell\right\}\right).
\end{equation}

\smallskip
\noindent\textup{(ii) Complete exponential Bell polynomial:}
\begin{equation}\label{eq:C-complete-Bell}
 C_m(r)=\frac1{(2m)!}\,
 \Gaus\!\left[
 Y_{2m}\bigl(1!h_1,2!h_2,\ldots,(2m)!h_{2m}\bigr)
 \right],
\end{equation}
where $Y_j$ is the $j$th complete exponential Bell polynomial.

\smallskip
\noindent\textup{(iii) Triangular recurrence:}
\begin{equation}\label{eq:P-recurrence}
 P_0=1,
 \qquad
 P_\nu=\frac1\nu\sum_{\ell=1}^{\nu}
 \ell h_\ell P_{\nu-\ell},
 \qquad C_m=\Gaus[P_{2m}].
\end{equation}

\smallskip
\noindent\textup{(iv) Integer-partition sum:}
if $\lambda=(1^{m_1}2^{m_2}\cdots)\vdash2m$ and
$K=\ell(\lambda)=\sum_{j\ge1}m_j$, then
\begin{equation}\label{eq:C-partition}
 C_m(r)=\sum_{\lambda\vdash2m}
 \frac{(-1)^{m+K}(2m+2K-1)!!}
      {[r(1+r)]^{m+K}}
 \prod_{j\ge1}
 \frac{\Touch_{j+2}(r)^{m_j}}
      {m_j!\,((j+2)!)^{m_j}}.
\end{equation}
Equivalently, summing over nonnegative integers
$k_1,\ldots,k_{2m}$ constrained by
$\sum_{\ell=1}^{2m}\ell k_\ell=2m$,
\begin{equation}\label{eq:C-multiindex}
 C_m(r)=
 \sum_{\sum \ell k_\ell=2m}
 \frac{(-1)^{m+K}(2m+2K-1)!!}
      {[r(1+r)]^{m+K}}
 \prod_{\ell=1}^{2m}
 \frac{\Touch_{\ell+2}(r)^{k_\ell}}
      {k_\ell!\,((\ell+2)!)^{k_\ell}},
 \quad K=\sum k_\ell.
\end{equation}
\end{theorem}

\begin{proof}
Equation \eqref{eq:C-gaussian} is the coefficient obtained by inserting
\eqref{eq:scaled-phase} into the central integral in the proof of
\cref{thm:principal}.  The standard generating identity
\[
 \exp\!\left(\sum_{j\ge1}x_j\frac{t^j}{j!}\right)
 =\sum_{\nu\ge0}Y_\nu(x_1,\ldots,x_\nu)\frac{t^\nu}{\nu!}
\]
gives \eqref{eq:C-complete-Bell}.  Differentiating
\eqref{eq:P-generating} with respect to $t$ and comparing coefficients gives
\eqref{eq:P-recurrence}.

For the partition formula, choose $k_\ell$ copies of $h_\ell t^\ell$ from the
exponential.  The constraint is $\sum\ell k_\ell=2m$.  If
$K=\sum k_\ell$, the resulting power of $u$ is
\[
 \sum_{\ell\ge1}(\ell+2)k_\ell=2m+2K.
\]
Consequently the power of $\ii$ is $(-1)^{m+K}$ and the Gaussian moment is
$(2m+2K-1)!!$.  Substitution of \eqref{eq:alpha} gives
\eqref{eq:C-multiindex}, which is the multiplicity notation for
\eqref{eq:C-partition}.
\end{proof}

\subsection{Rational structure and degree bounds}

\begin{theorem}[Denominator theorem]\label{thm:denominator}
For every $m\ge0$ there is a polynomial $P_m(r)\in\mathbb Q[r]$ such that
\begin{equation}\label{eq:C-rational-structure}
 C_m(r)=\frac{P_m(r)}{r^m(1+r)^{3m}},
 \qquad \deg P_m\le4m.
\end{equation}
In particular, $C_m$ is regular for $r>0$ and has a finite limit as
$r\to\infty$.
\end{theorem}

\begin{proof}
Every $\Touch_j(r)$ with $j\ge1$ contains a factor $r$.  In a summand of
\eqref{eq:C-partition} with length $K$, these factors cancel $r^K$ from
$[r(1+r)]^{m+K}$, leaving a denominator
$r^m(1+r)^{m+K}$.  Since a partition of $2m$ has $K\le2m$, a common
denominator is $r^m(1+r)^{3m}$.

After removing the factor $r$ from each $\Touch_{j+2}(r)$, the numerator degree
in that summand is at most
\[
 \sum_{j\ge1}(j+1)m_j=2m+K.
\]
Multiplication by $(1+r)^{2m-K}$ to reach the common denominator raises this by
at most $2m-K$, giving degree at most $4m$.
\end{proof}

\begin{corollary}[Finite algorithm and arithmetic]\label{cor:finite-algorithm}
For a fixed $m$, computing $C_m$ requires only the Touchard polynomials through
index $2m+2$ and the integer partitions of $2m$.  All arithmetic can be carried
out exactly over $\mathbb Q(r)$.
\end{corollary}

\begin{remark}
The number of summands in \eqref{eq:C-partition} is the partition number
$p(2m)$, rather than the much larger number of terms obtained by naively
expanding repeated derivatives.  This is the main practical advantage of the
partition organization.
\end{remark}

\section{Explicit principal coefficients}

The first coefficient follows from the two partitions $2$ and $1+1$:
\begin{equation}\label{eq:C1-alpha}
 C_1=3\alpha_4-\frac{15}{2}\alpha_3^2.
\end{equation}
After substitution of the Touchard polynomials, the first five coefficients
are as follows.
\begin{align}
 C_0(r)&=1,\label{eq:C0-explicit}\\[0.25em]
 C_1(r)&=-\frac{2r^4+9r^3+16r^2+6r+2}
 {24r(1+r)^3},\label{eq:C1-explicit}\\[0.25em]
 C_2(r)&=\frac{
 4r^8-12r^7-143r^6-384r^5-588r^4-636r^3
 +100r^2+24r+4}
 {1152r^2(1+r)^6},\label{eq:C2-explicit}
\end{align}
\begin{align}
 C_3(r)&=\frac{1}{414720r^3(1+r)^9}
 \Bigl(
 1112r^{12}+12708r^{11}+56082r^{10}+118683r^9\notag\\
 &\hspace{2.6em}{}+97176r^8-74394r^7-259046r^6-456444r^5
 -635424r^4\notag\\
 &\hspace{2.6em}{}+115308r^3+39432r^2+10008r+1112
 \Bigr),\label{eq:C3-explicit}
\end{align}
\begin{align}
 C_4(r)&=-\frac{1}{39813120r^4(1+r)^{12}}
 \Bigl(
 9136r^{16}-357408r^{15}-5092600r^{14}-29432904r^{13}\notag\\
 &\hspace{2.4em}{}-98499653r^{12}-212369184r^{11}
 -302729704r^{10}-277448328r^9\notag\\
 &\hspace{2.4em}{}-156972248r^8-52300272r^7+110053352r^6
 +302787360r^5\notag\\
 &\hspace{2.4em}{}+198352r^4+2545056r^3+647456r^2+109632r+9136
 \Bigr).\label{eq:C4-explicit}
\end{align}
The apparently small coefficient $198352$ of $r^4$ in
\eqref{eq:C4-explicit} is not a typographical omission; it follows from the
exact partition sum and was checked independently by recurrence and direct
series expansion.

\begin{resultbox}[title=First practical truncations,fonttitle=\bfseries]
With $q=r/n$, the one-, two-, and three-correction approximations are
\begin{align*}
 \Bell_n&=L_n(r)\{1+\bigO(q)\},\\
 \Bell_n&=L_n(r)\{1+C_1(r)q+\bigO(q^2)\},\\
 \Bell_n&=L_n(r)\{1+C_1(r)q+C_2(r)q^2+\bigO(q^3)\}.
\end{align*}
For numerical work, the first neglected term is a useful local diagnostic, but
it is not by itself a rigorous global error bound once the series is taken near
optimal truncation.
\end{resultbox}

\section{The Stirling interface}

\subsection{The limiting coefficients}

Let $\widehat s_m$ denote the reciprocal Stirling coefficients, defined by
\begin{equation}\label{eq:reciprocal-stirling}
 \frac1{\Gamma(x+1)}
 \sim\frac{\ee^x x^{-x-1/2}}{\sqrt{2\pi}}
 \sum_{m\ge0}\widehat s_m x^{-m}.
\end{equation}
Their ordinary generating series is the formal exponential
\begin{equation}\label{eq:reciprocal-stirling-gen}
 \sum_{m\ge0}\widehat s_m z^m
 =\exp\!\left(
 -\sum_{\nu\ge1}
 \frac{B_{2\nu}}{2\nu(2\nu-1)}z^{2\nu-1}
 \right),
\end{equation}
where $B_{2\nu}$ are Bernoulli numbers.

\begin{theorem}[Reciprocal-Stirling limit]\label{thm:stirling-limit}
For every fixed $m\ge0$,
\begin{equation}\label{eq:C-limit}
 \lim_{r\to\infty}C_m(r)=\widehat s_m.
\end{equation}
In particular,
\begin{equation}\label{eq:C-limits-list}
 1,\quad -\frac1{12},\quad \frac1{288},\quad
 \frac{139}{51840},\quad -\frac{571}{2488320},\ldots
\end{equation}
are the limits of $C_0,C_1,C_2,C_3,C_4,\ldots$.
\end{theorem}

\begin{proof}
For fixed $j$,
\[
 \alpha_j(r)=\frac{\Touch_j(r)}{j![r(1+r)]^{j/2}}
 \longrightarrow\frac1{j!}.
\]
Hence \eqref{eq:C-gaussian} tends coefficientwise to
\begin{equation}\label{eq:limit-gaussian}
 \Gaus\!\left([t^{2m}]
 \exp\!\left\{\sum_{j\ge3}\frac{(\ii u)^j}{j!}t^{j-2}\right\}\right).
\end{equation}
To identify these numbers, use the exact diagonal integral
\begin{equation}\label{eq:poisson-diagonal}
 \frac1{2\pi}\int_{-\pi}^{\pi}
 \exp\!\bigl(x\ee^{\ii\theta}-x-\ii x\theta\bigr)\dd\theta
 =\frac{\ee^{-x}x^x}{\Gamma(x+1)}
\end{equation}
for positive integer $x$, and apply the same Gaussian scaling
$\theta=u/\sqrt{x}$.  Stirling's expansion of the right-hand side identifies
\eqref{eq:limit-gaussian} with $\widehat s_m$.  Since both sides are formal
coefficient identities, the conclusion holds for every $m$.
\end{proof}

\begin{remark}
Hwang proved a broader formal identity showing that two apparently different
saddle coefficient systems both reduce to the same Stirling coefficients.
Here only the specialization needed for the Bell circle saddle is used.
\end{remark}

\subsection{Factorial-free all-order expansion}

Let the ordinary Stirling coefficients $s_j$ be defined by
\begin{equation}\label{eq:stirling-series}
 n!\sim\sqrt{2\pi n}\left(\frac ne\right)^n
 \sum_{j\ge0}s_jn^{-j},
\end{equation}
with
\begin{equation}\label{eq:stirling-gen}
 \sum_{j\ge0}s_jz^j
 =\exp\!\left(
 \sum_{\nu\ge1}
 \frac{B_{2\nu}}{2\nu(2\nu-1)}z^{2\nu-1}
 \right).
\end{equation}
Thus $s_0=1$, $s_1=1/12$, $s_2=1/288$,
$s_3=-139/51840$, and so on.

\begin{theorem}[Factorial-free form]\label{thm:factorial-free}
Put
\begin{equation}\label{eq:Dm}
 D_m(r):=r^mC_m(r),
 \qquad
 \cE_\ell(r):=\sum_{m=0}^{\ell}s_{\ell-m}D_m(r).
\end{equation}
Then
\begin{equation}\label{eq:factorial-free}
 \Bell_n\sim
 \frac{\exp\{nr-n+n/r-1\}}{\sqrt{1+r}}
 \sum_{\ell\ge0}\frac{\cE_\ell(r)}{n^\ell},
 \qquad r=\W_0(n).
\end{equation}
Although written in powers of $1/n$, this is still an expansion in the
asymptotic scale $(r/n)^\ell$, because $\cE_\ell(r)=\bigO(r^\ell)$.
\end{theorem}

\begin{proof}
Insert \eqref{eq:stirling-series} into
\eqref{eq:principal-theorem} and use
$\log(n/r)=r$ and $\ee^r=n/r$.  Cauchy multiplication gives
\eqref{eq:Dm}.  The growth statement follows from
\cref{thm:stirling-limit}.
\end{proof}

The first coefficients in this normalization are
\begin{align}
 \cE_0(r)&=1,\label{eq:E0}\\
 \cE_1(r)&=-\frac{r^2(2r^2+7r+10)}{24(1+r)^3},\label{eq:E1}\\
 \cE_2(r)&=\frac{r^3(4r^5-20r^4-199r^3-556r^2-860r-864)}
 {1152(1+r)^6},\label{eq:E2}\\
 \cE_3(r)&=\frac{r^4}{414720(1+r)^9}
 \Bigl(1112r^8+12828r^7+55962r^6+111301r^5\notag\\
 &\hspace{7em}{}+55818r^4-196476r^3-492584r^2
 -749376r-867456\Bigr).\label{eq:E3}
\end{align}
The negative sign in \eqref{eq:E1} is important: it follows from a cancellation
between the $+1/(12n)$ Stirling correction and the larger saddle correction
$rC_1(r)/n$.

\section{The logarithmic transseries}

Multiplicative asymptotic series are often best manipulated after taking a
formal logarithm.  Define
\begin{equation}\label{eq:F-Lambda}
 F(q;r):=1+\sum_{m\ge1}C_m(r)q^m,
 \qquad
 \log F(q;r)=\sum_{m\ge1}\Lambda_m(r)q^m.
\end{equation}

\begin{proposition}[Cumulant recurrence]\label{prop:lambda-recurrence}
The coefficients $\Lambda_m$ are determined by
\begin{equation}\label{eq:lambda-recurrence}
 \Lambda_m(r)=C_m(r)-\frac1m
 \sum_{j=1}^{m-1}j\Lambda_j(r)C_{m-j}(r),
 \qquad m\ge1.
\end{equation}
Thus
\begin{align*}
 \Lambda_1&=C_1,\\
 \Lambda_2&=C_2-\frac12C_1^2,\\
 \Lambda_3&=C_3-C_1C_2+\frac13C_1^3.
\end{align*}
\end{proposition}

\begin{proof}
Differentiate $F=\exp(\log F)$, use $F'=F(\log F)'$, and compare the
coefficient of $q^{m-1}$.
\end{proof}

The first two nonzero cumulants simplify to
\begin{align}
 \Lambda_1(r)&=-\frac{2r^4+9r^3+16r^2+6r+2}
 {24r(1+r)^3},\label{eq:Lambda1}\\
 \Lambda_2(r)&=-\frac{r(2r^4+12r^3+29r^2+40r+36)}
 {48(1+r)^6}.\label{eq:Lambda2}
\end{align}

\begin{theorem}[All-order logarithm]\label{thm:logarithm}
For fixed truncation order, the logarithm of the Bell number has the expansion
\begin{align}
 \log\Bell_n
 &\sim n\left(r-1+\frac1r\right)-1-\frac12\log(1+r)\notag\\
 &\quad+\sum_{\nu\ge1}
 \frac{B_{2\nu}}{2\nu(2\nu-1)n^{2\nu-1}}
 +\sum_{m\ge1}\Lambda_m(r)\left(\frac rn\right)^m,
 \qquad r=\W_0(n).
 \label{eq:log-all-order}
\end{align}
Equivalently,
\begin{equation}\label{eq:log-Hj}
 \log\Bell_n\sim nA(r)+B(r)+\sum_{j\ge1}\frac{H_j(r)}{n^j},
\end{equation}
where
\begin{align}
 A(r)&:=r-1+r^{-1},\label{eq:A-r}\\
 B(r)&:=-1-\tfrac12\log(1+r),\label{eq:B-r}\\
 H_j(r)&:=r^j\Lambda_j(r)
 +\ind_{\{j\ \mathrm{odd}\}}\frac{B_{j+1}}{j(j+1)}.
 \label{eq:Hj}
\end{align}
Every $H_j$ is an explicitly computable rational function of $r$.
\end{theorem}

\begin{proof}
Take the logarithm of \eqref{eq:principal-theorem}, use the Stirling expansion
for $\log n!$, and simplify with $\log(n/r)=r$ and $\ee^r=n/r$.
The formal logarithm of the correction factor is
\eqref{eq:F-Lambda}.  Formula \eqref{eq:Hj} only regroups terms carrying the
same power of $1/n$.
\end{proof}

\begin{remark}[Residual certificates]
Suppose $F_M=1+\sum_{m=1}^{M-1}C_mq^m$.  Formal division verifies
\[
 \exp\!\left(\sum_{m=1}^{M-1}\Lambda_mq^m\right)-F_M
 =\bigO(q^M).
\]
This identity is a simple symbolic certificate that the multiplicative and
logarithmic truncations agree through the requested order.
\end{remark}

\section{Lambert--\texorpdfstring{$W$}{W} polynomial--logarithmic re-expansion}

\subsection{Every branch}

Let $z$ tend to infinity away from the branch cuts and put
\begin{equation}\label{eq:xi-k}
 \xi_k:=\Log z+2\pi\ii k.
\end{equation}
Let $\Sone{a}{b}$ denote an unsigned Stirling number of the first
kind.  Define the Lambert polynomials
\begin{equation}\label{eq:Lambert-polynomials}
 \cL_a(y):=\sum_{b=1}^{a}(-1)^{a-b}
 \Sone{a}{a-b+1}\frac{y^b}{b!},
 \qquad a\ge1.
\end{equation}
The standard branch expansion can then be written compactly as
\begin{equation}\label{eq:W-all-branches}
 \W_k(z)\sim \xi_k-\log\xi_k
 +\sum_{a\ge1}\frac{\cL_a(\log\xi_k)}{\xi_k^a}.
\end{equation}
This formula supplies every coefficient using only Stirling numbers of the
first kind.  The first polynomials are
\begin{align*}
 \cL_1(y)&=y,\\
 \cL_2(y)&=\frac{y^2}{2}-y,\\
 \cL_3(y)&=\frac{y^3}{3}-\frac32y^2+y,\\
 \cL_4(y)&=\frac{y^4}{4}-\frac{11}{6}y^3+3y^2-y,\\
 \cL_5(y)&=\frac{y^5}{5}-\frac{25}{12}y^4
 +\frac{35}{6}y^3-5y^2+y.
\end{align*}

For the principal real branch, write
\begin{equation}\label{eq:L-ell}
 L:=\log n,\qquad \ell:=\log L.
\end{equation}
Then
\begin{align}
 r=\W_0(n)
 &\sim L-\ell+\frac{\ell}{L}
 +\frac{\ell(\ell-2)}{2L^2}
 +\frac{\ell(6-9\ell+2\ell^2)}{6L^3}\notag\\
 &\quad+\frac{\ell(-12+36\ell-22\ell^2+3\ell^3)}{12L^4}
 +\cdots.\label{eq:W-principal-expanded}
\end{align}

\subsection{Canonical nested-log form of \texorpdfstring{$\log\Bell_n$}{log Bn}}

Substitution of \eqref{eq:W-principal-expanded} into
\eqref{eq:log-all-order} produces the polynomial--logarithmic transseries
without any further inversion.  The first inverse-log layers, before the
algebraic $1/n$ layer, are
\begin{align}
 \log\Bell_n
 &\sim n\Biggl[
 L-\ell-1+\frac{\ell+1}{L}
 +\frac{\ell^2}{2L^2}
 +\frac{\ell^3/3-\ell^2/2}{L^3}\notag\\
 &\hspace{5.2em}{}
 +\frac{\ell^4/4-5\ell^3/6+\ell^2/2}{L^4}
 +\cdots\Biggr]\notag\\
 &\quad-\frac12\log L-1
 +\frac{\ell-1}{2L}
 +\frac{\ell^2-4\ell+1}{4L^2}\notag\\
 &\quad+\frac{2\ell^3-15\ell^2+18\ell-2}{12L^3}
 +\cdots
 +\sum_{j\ge1}\frac{H_j(\W_0(n))}{n^j}.
 \label{eq:log-nested}
\end{align}
The final sum in \eqref{eq:log-nested} is itself expanded by
\eqref{eq:W-principal-expanded}.  It starts at the much smaller scale
$\log n/n$.

\begin{algorithm}[Mechanical nested-log expansion]\label{alg:nested-log}
To compute \eqref{eq:log-nested} through any prescribed monomial order:
\begin{enumerate}
 \item generate $\cL_a$ from \eqref{eq:Lambert-polynomials};
 \item truncate \eqref{eq:W-all-branches} at the required inverse-log order;
 \item generate $C_m$ from \eqref{eq:C-partition}, then $\Lambda_m$ from
       \eqref{eq:lambda-recurrence};
 \item substitute the truncated $W$ series into $A,B,H_j$ in
       \eqref{eq:log-Hj};
 \item expand in the ordered monomial family
       $n^{-j}L^{-a}\ell^b$, retaining the chosen finite lower set.
\end{enumerate}
Because every step is a finite polynomial or rational operation at a fixed
order, the procedure terminates and gives a residual certificate by direct
substitution.
\end{algorithm}

\section{The full Lambert-branch saddle transseries}

The principal circle expansion cannot see exponentially small saddles one at a
time.  For that purpose we return to the Hankel integral \eqref{eq:hankel}.

\subsection{Local data at an arbitrary saddle}

At $w=w_k=\W_k(N)$,
\begin{equation}\label{eq:psi-at-saddle}
 \psi_N(w)=\frac1w-\Log w,
 \qquad
 \psi_N''(w)=\frac{1+w}{w^2},
\end{equation}
and for $j\ge2$,
\begin{equation}\label{eq:psi-derivs}
 \psi_N^{(j)}(w)=\frac1w+(-1)^j\frac{(j-1)!}{w^j}.
\end{equation}
Use the local steepest coordinate
\begin{equation}\label{eq:vertical-coordinate}
 t=w+\frac{\ii w}{\sqrt{N(1+w)}}u.
\end{equation}
Define
\begin{equation}\label{eq:beta}
 \beta_j(w):=
 \frac{w^j\psi_N^{(j)}(w)}{j!(1+w)^{j/2}}
 =\frac{w^{j-1}+(-1)^j(j-1)!}{j!(1+w)^{j/2}},
 \qquad j\ge3.
\end{equation}
Then the local exponent is
\begin{equation}\label{eq:vertical-scaled-phase}
 N\{\psi_N(t)-\psi_N(w)\}
 =-\frac{u^2}{2}
 +\sum_{j\ge3}\beta_j(w)(\ii u)^jN^{-(j-2)/2}.
\end{equation}

\subsection{All coefficients in one thimble sector}

For $\ell\ge1$, set
\begin{equation}\label{eq:tilde-h}
 \widetilde h_\ell(u;w):=\beta_{\ell+2}(w)(\ii u)^{\ell+2}.
\end{equation}
Define $\widetilde P_\nu$ by the same triangular recurrence as before:
\begin{equation}\label{eq:tilde-P-rec}
 \widetilde P_0=1,
 \qquad
 \widetilde P_\nu=\frac1\nu\sum_{\ell=1}^{\nu}
 \ell\widetilde h_\ell\widetilde P_{\nu-\ell}.
\end{equation}

\begin{theorem}[Arbitrary Lambert-sector expansion]\label{thm:sector}
For a fixed saddle branch $k$ whose thimble contributes, let
$w_k=\W_k(N)$.  Its oriented contribution has the formal Poincar\'e expansion
\begin{equation}\label{eq:Jk-expansion}
 \cJ_k(N)\sim
 \frac{\Gamma(N)\exp(-1+N/w_k)}
 {w_k^{N-1}\sqrt{2\pi N(1+w_k)}}
 \sum_{m\ge0}\frac{\widetilde C_m(w_k)}{N^m},
\end{equation}
where the square root is fixed by the thimble orientation and
\begin{equation}\label{eq:Ctilde-gauss}
 \widetilde C_m(w)=\Gaus[\widetilde P_{2m}(u;w)].
\end{equation}
Equivalently,
\begin{equation}\label{eq:Ctilde-partition}
 \boxed{
 \widetilde C_m(w)=
 \sum_{\lambda\vdash2m}
 (-1)^{m+K}(2m+2K-1)!!
 \prod_{j\ge1}
 \frac{\beta_{j+2}(w)^{m_j}}{m_j!},
 \qquad K=\ell(\lambda).}
\end{equation}
Moreover,
\begin{equation}\label{eq:Ctilde-rational}
 \widetilde C_m(w)=
 \frac{\widetilde P_m(w)}{(1+w)^{3m}},
 \qquad
 \widetilde P_m\in\mathbb Q[w],
 \quad \deg\widetilde P_m\le4m.
\end{equation}
\end{theorem}

\begin{proof}
Insert \eqref{eq:vertical-coordinate} into the thimble integral and use
\eqref{eq:vertical-scaled-phase}.  The Gaussian coefficient extraction and
partition expansion are identical to the proof of
\cref{thm:coefficient-calculus}, with $q^{1/2}$ replaced by $N^{-1/2}$ and
$\alpha_j$ replaced by $\beta_j$.  Each factor $\beta_{j+2}$ has denominator
$(1+w)^{(j+2)/2}$.  A partition of $2m$ of length $K$ therefore contributes
$(1+w)^{-(m+K)}$; since $K\le2m$, the common denominator is
$(1+w)^{3m}$.  The numerator-degree estimate is the same as in
\cref{thm:denominator}.
\end{proof}

The first coefficients are
\begin{align}
 \widetilde C_0(w)&=1,\label{eq:Ct0}\\
 \widetilde C_1(w)&=-\frac{2w^4-3w^3-20w^2-18w+2}
 {24(1+w)^3},\label{eq:Ct1}\\
 \widetilde C_2(w)&=\frac{
 4w^8-156w^7-695w^6-696w^5+1092w^4+2916w^3
 +1972w^2-72w+4}
 {1152(1+w)^6}.\label{eq:Ct2}
\end{align}
In Paris's notation these are
$\widetilde C_m=c_{2m}(1/2)_m$; thus
\eqref{eq:Ct1}--\eqref{eq:Ct2} agree with the published low-order saddle
coefficients, while \eqref{eq:Ctilde-partition} supplies all higher orders.

\subsection{Thimble selection and Stokes multipliers}

For positive real argument, the saddles occur in conjugate pairs.  Deforming
the original Hankel loop into steepest-descent paths gives the exact contour
decomposition
\begin{equation}\label{eq:exact-thimble-decomp}
 \Bell_n=\cJ_0(N)+2\RePart\sum_{k=1}^{K(N)}\cJ_k(N),
 \qquad N=n+1,
\end{equation}
away from a Stokes threshold.  The integer $K(N)$ is determined by thimble
connectivity.  It increases when the positive-axis parameter $N$ crosses a
Stokes boundary.  The first boundaries affecting the positive axis are
\begin{center}
\begin{tabular}{@{}ccc@{}}
\toprule
range of $N$ & active positive indices & total saddle count\\
\midrule
$0<N<6.87877\ldots$ & $k=1$ & $3$\\
$6.87877\ldots<N<13.56411\ldots$ & $k=1,2$ & $5$\\
$13.56411\ldots<N<20.13877\ldots$ & $k=1,2,3$ & $7$\\
$20.13877\ldots<N<26.43594\ldots$ & $k=1,2,3,4$ & $9$\\
\bottomrule
\end{tabular}
\end{center}
The omitted boundaries continue by the same adjacent-thimble Stokes test.  A
computational definition is: a new pair is switched on when adjacent saddle
actions have equal imaginary part,
\begin{equation}\label{eq:stokes-condition}
 \ImPart\{\psi_N(w_j)-\psi_N(w_{j+1})\}=0,
\end{equation}
and the corresponding steepest path has the required connectivity to the
original loop.  The connectivity clause is essential; equality of phases
alone does not determine an intersection number.

It is convenient to encode this by Stokes multipliers
$\sigma_k(N)\in\{0,1\}$ away from a transition:
\begin{equation}\label{eq:formal-full-transseries}
 \boxed{
 \Bell_n\sim
 \cJ_0^{\mathrm{formal}}(N)
 +2\RePart\sum_{k\ge1}\sigma_k(N)
 \cJ_k^{\mathrm{formal}}(N),}
\end{equation}
where $\cJ_k^{\mathrm{formal}}$ denotes \eqref{eq:Jk-expansion}.  At every
fixed $N$ only finitely many $\sigma_k$ are nonzero.  A median convention gives
$\sigma_k=1/2$ for a sector exactly on its Stokes line; lateral contours give
the two one-sided values.

\begin{remark}[Meaning of the global formula]
Equation \eqref{eq:exact-thimble-decomp} is an exact contour statement.
Equation \eqref{eq:formal-full-transseries} attaches a Poincar\'e series to each
exact thimble.  The statement is sectorwise, not a claim that all active
series may be truncated at one common order with an error estimate uniform in
$k$ as $K(N)\to\infty$.
\end{remark}

\subsection{Relation between the two principal normalizations}

The $k=0$ sector in \eqref{eq:Jk-expansion} uses
$w_0=\W_0(n+1)$ and a $1/(n+1)$ expansion.  The circle result uses
$r=\W_0(n)$ and $q=r/n$.  Both are expansions of the same exact quantity, but
their coefficient functions are not identical term by term.  They are related
by the formal shift
\begin{equation}\label{eq:w-r-shift}
 w_0=r+\frac{r}{n(1+r)}+\bigO\!\left(\frac{r^2}{n^2}\right)
\end{equation}
and subsequent re-expansion.  Keeping the two forms separate makes each local
calculus simpler and prevents accidental mixing of $n$ and $n+1$.

\section{Singulants and beyond-all-orders ordering}

\subsection{Exact singulant}

Let $R:=w_0=\W_0(N)$.  Compare the leading term of the $k$th sector with the
principal one.  Define
\begin{equation}\label{eq:singulant}
 \cA_k(N):=N\left\{
 \Log\frac{w_k}{R}+\frac1R-\frac1{w_k}
 \right\}.
\end{equation}
Then
\begin{equation}\label{eq:sector-ratio}
 \frac{\cJ_k^{(0)}(N)}{\cJ_0^{(0)}(N)}
 =\exp\{-\cA_k(N)\}\,
 \frac{w_k}{R}\sqrt{\frac{1+R}{1+w_k}}.
\end{equation}
The logarithms in \eqref{eq:singulant} are chosen by continuation from the
corresponding $W$ branches, consistently with the thimble orientation.

\subsection{Fixed-branch expansion}

Set
\begin{equation}\label{eq:delta-def}
 a_k:=2\pi\ii k,
 \qquad \delta_k:=w_k-R.
\end{equation}
Subtracting $R+\log R=\log N$ from
$w_k+\Log w_k=\log N+2\pi\ii k$ gives the exact implicit equation
\begin{equation}\label{eq:delta-implicit}
 \delta_k+\log\!\left(1+\frac{\delta_k}{R}\right)=a_k.
\end{equation}
Its first terms are
\begin{align}
 \delta_k
 &=a_k-\frac{a_k}{R}
 +\frac{a_k(a_k+2)}{2R^2}
 -\frac{a_k(2a_k^2+9a_k+6)}{6R^3}\notag\\
 &\quad+\frac{a_k(3a_k^3+22a_k^2+36a_k+12)}{12R^4}
 +\bigO_k(R^{-5}).\label{eq:delta-expanded}
\end{align}
All coefficients are generated by the finite recurrence below.  If
$\delta(a,x)=\sum_{m\ge0}d_m(a)x^m$ with $x=R^{-1}$ and $d_0(a)=a$, then
\begin{equation}\label{eq:delta-recurrence}
 d_m(a)=-\sum_{j=1}^{m}\frac{(-1)^{j+1}}{j}
 [x^{m-j}]\,\delta(a,x)^j,
 \qquad m\ge1.
\end{equation}

Substitution into \eqref{eq:singulant} yields
\begin{align}
 \frac{\cA_k(N)}{N}
 &=\frac{a_k}{R}-\frac{a_k^2}{2R^2}
 +\frac{a_k^3/3+a_k^2/2}{R^3}\notag\\
 &\quad-\frac{a_k^4/4+5a_k^3/6+a_k^2/2}{R^4}
 +\bigO_k(R^{-5}).\label{eq:singulant-expanded}
\end{align}
In particular,
\begin{equation}\label{eq:singulant-real}
 \RePart\cA_k(N)=
 \frac{2\pi^2k^2N}{R^2}
 \left(1-\frac1R+\bigO_k(R^{-2})\right).
\end{equation}

\begin{corollary}[Beyond every algebraic order]\label{cor:beyond-all-orders}
For each fixed $k\ne0$ and every fixed $M>0$,
\begin{equation}\label{eq:beyond-all-orders}
 \frac{\cJ_k(N)}{\cJ_0(N)}
 =\smallo\!\left(\left(\frac{R}{N}\right)^M\right)
 \qquad(N\to\infty),
\end{equation}
sectorwise after fixed-order truncation.  More explicitly, the leading
magnitude is
\begin{equation}\label{eq:sector-magnitude}
 \left|\frac{\cJ_k^{(0)}}{\cJ_0^{(0)}}\right|
 =\exp\!\left\{-\frac{2\pi^2k^2N}{R^2}
 \left(1+\bigO_k(R^{-1})\right)\right\}\bigO_k(1).
\end{equation}
\end{corollary}

\begin{proof}
Equation \eqref{eq:singulant-real} shows that the logarithm of the sector ratio
is $-\Theta_k(N/\log^2N)$, whereas the logarithm of any fixed power of $R/N$
is only $-\Theta_M(\log N)$.  Therefore the exponential sector is smaller.
\end{proof}

\begin{remark}[Oscillations]
The imaginary part of $\cA_k$ begins with
$2\pi kN/R-8\pi^3k^3N/(3R^3)+\cdots$.  The conjugate pair
$k,-k$ consequently contributes an exponentially damped cosine with a rapidly
varying phase.  The exact phase also includes the argument of the prefactor in
\eqref{eq:sector-ratio} and the phase of the truncated coefficient series.
\end{remark}

\section{Executable coefficient generation}

The following compact SymPy program implements the partition formula exactly.
It generates $C_m$, verifies the denominator theorem, and also generates the
arbitrary-thimble coefficients $\widetilde C_m$.

\begin{lstlisting}[style=code,caption={Exact principal and thimble coefficient generation.}]
import sympy as sp

r, w = sp.symbols("r w")


def touchard(j, x):
    return sp.expand(sum(
        sp.functions.combinatorial.numbers.stirling(j, a, kind=2) * x**a
        for a in range(j + 1)
    ))


def multiplicity_partitions(total):
    """Yield dictionaries {part: multiplicity} for partitions of total."""
    def rec(rem, part, data):
        if rem == 0:
            yield dict(data)
            return
        if part == 0:
            return
        for count in range(rem // part, -1, -1):
            if count:
                data[part] = count
            else:
                data.pop(part, None)
            yield from rec(rem - count * part, part - 1, data)
        data.pop(part, None)
    yield from rec(total, total, {})


def principal_C(m):
    if m == 0:
        return sp.Integer(1)
    ans = 0
    for mult in multiplicity_partitions(2 * m):
        K = sum(mult.values())
        term = (-1)**(m + K) * sp.factorial2(2 * m + 2 * K - 1)
        term /= (r * (1 + r))**(m + K)
        for ell, count in mult.items():
            term *= touchard(ell + 2, r)**count
            term /= sp.factorial(count) * sp.factorial(ell + 2)**count
        ans += term
    return sp.factor(sp.cancel(ans))


def beta(j):
    return (w**(j - 1) + (-1)**j * sp.factorial(j - 1)) / (
        sp.factorial(j) * (1 + w)**sp.Rational(j, 2)
    )


def thimble_C(m):
    if m == 0:
        return sp.Integer(1)
    ans = 0
    for mult in multiplicity_partitions(2 * m):
        K = sum(mult.values())
        term = (-1)**(m + K) * sp.factorial2(2 * m + 2 * K - 1)
        for ell, count in mult.items():
            term *= beta(ell + 2)**count / sp.factorial(count)
        ans += term
    return sp.factor(sp.cancel(ans))


for m in range(5):
    Cm = principal_C(m)
    print(m, Cm)
    print("denominator check:",
          sp.cancel(Cm * r**m * (1 + r)**(3*m)))

for m in range(3):
    print("thimble", m, thimble_C(m))
\end{lstlisting}

\subsection{Stable numerical evaluation}

For large $n$, do not form $n!$ or $\Bell_n$ directly in floating-point
arithmetic.  Compute
\begin{equation}\label{eq:log-prefactor}
 \log L_n(r)=\log\Gamma(n+1)+\ee^r-1-n\log r
 -\frac12\log\{2\pi n(1+r)\}
\end{equation}
and evaluate the correction polynomial in $q$ by Horner's rule.  For a signed
relative approximation, return
\begin{equation}\label{eq:numeric-eval}
 \widehat{\Bell}_n^{(M)}=
 \exp\{\log L_n(r)\}
 \sum_{m=0}^{M-1}C_m(r)q^m.
\end{equation}
When only $\log\Bell_n$ is required, use \eqref{eq:log-all-order} directly; this
is both more stable and more naturally compatible with transseries
composition.

\section{Numerical validation}

Let $\widehat{\Bell}_n^{(M)}$ denote \eqref{eq:numeric-eval}, with $M$ terms
$C_0,\ldots,C_{M-1}$, and define the signed relative error
\begin{equation}\label{eq:rel-error}
 \varepsilon_{n,M}:=\frac{\widehat{\Bell}_n^{(M)}}{\Bell_n}-1.
\end{equation}
The exact Bell numbers in the following table were generated by the integer
recurrence
\begin{equation}\label{eq:Bell-recurrence}
 \Bell_0=1,
 \qquad
 \Bell_n=\sum_{j=0}^{n-1}\binom{n-1}{j}\Bell_j.
\end{equation}

\begin{center}
\scriptsize
\setlength{\tabcolsep}{4.2pt}
\begin{tabular}{@{}r r r r r r@{}}
\toprule
$n$ & $M=1$ & $M=2$ & $M=3$ & $M=4$ & $M=5$\\
\midrule
10
& $ 2.6771612450\,10^{-2}$
& $ 3.8253602751\,10^{-4}$
& $-1.2965341197\,10^{-5}$
& $-6.2268397638\,10^{-6}$
& $-1.1599280809\,10^{-6}$\\
20
& $ 1.5326827723\,10^{-2}$
& $ 1.1293363992\,10^{-4}$
& $-3.8237416551\,10^{-6}$
& $-7.5548760506\,10^{-7}$
& $-7.1302490456\,10^{-8}$\\
50
& $ 7.2381177878\,10^{-3}$
& $ 2.1197164537\,10^{-5}$
& $-5.8160606659\,10^{-7}$
& $-4.1043863046\,10^{-8}$
& $-1.3992028317\,10^{-9}$\\
100
& $ 4.0615707492\,10^{-3}$
& $ 5.7437480134\,10^{-6}$
& $-1.2399465806\,10^{-7}$
& $-4.2053807468\,10^{-9}$
& $-6.2657057618\,10^{-11}$\\
200
& $ 2.2599385476\,10^{-3}$
& $ 1.5057676926\,10^{-6}$
& $-2.4589871384\,10^{-8}$
& $-4.0769478428\,10^{-10}$
& $-2.4569359891\,10^{-12}$\\
500
& $ 1.0289535458\,10^{-3}$
& $ 2.4364162415\,10^{-7}$
& $-2.6649303424\,10^{-9}$
& $-1.7364703016\,10^{-11}$
& $-2.4871987030\,10^{-14}$\\
\bottomrule
\end{tabular}
\end{center}

Several features are visible.  The leading approximation is already within a
few percent at $n=10$.  The $C_1$ correction improves it by roughly two orders
of magnitude.  The signs need not alternate monotonically at small $n$, but
fixed-order accuracy rapidly approaches the predicted $q^M$ scale.  The
five-term error at $n=500$ is at the level of $2.5\times10^{-14}$ in the
working precision used for the calculation.

\section{Proof-audit and formalization map}

This section records the logical dependencies in a form suitable for later
formalization.

\subsection{Analytic kernel}

A proof assistant development would require the following analytic statements.
\begin{enumerate}
 \item Cauchy's coefficient formula specialized to the circle
       $z=r\ee^{\ii\theta}$, giving \eqref{eq:circle}.
 \item Existence and uniqueness of $r>0$ satisfying $r\ee^r=n$.
 \item The derivative identity \eqref{eq:derivative-touchard} and the exact
       phase expansion \eqref{eq:phase-taylor}.
 \item Uniform normalized derivative bounds such as
       \eqref{eq:alpha-bound}.
 \item The local quadratic decay \eqref{eq:local-decay} and the global bound
       \eqref{eq:global-real-bound}.
 \item A finite Taylor theorem under a Gaussian integral, with a remainder
       uniform in $r$ once $r\ge1$.
\end{enumerate}
These yield \cref{thm:principal} without invoking any unformalized general
``saddle-point method'' black box.

\subsection{Algebraic kernel}

The coefficient calculus reduces to finite combinatorics:
\begin{enumerate}
 \item exponential Bell polynomial generating identities;
 \item the recurrence \eqref{eq:P-recurrence};
 \item Gaussian moments \eqref{eq:gaussian-moments};
 \item the bijection between constrained multiplicity vectors and partitions;
 \item polynomial degree and denominator bookkeeping.
\end{enumerate}
This part is particularly well suited to exact symbolic verification.

\subsection{Transseries kernel}

The global extension additionally needs:
\begin{enumerate}
 \item the Hankel representation \eqref{eq:hankel};
 \item the branch equation and conjugacy properties of $\W_k$;
 \item oriented thimble deformation and intersection numbers;
 \item local saddle expansion on each thimble;
 \item a formal algebra for the branch expansion \eqref{eq:W-all-branches}
       and for composition in the monomial family
       \[n^{-j}L^{-a}\ell^b\ee^{-\cA_k}.\]
\end{enumerate}
The present article separates these analytic, algebraic, and transseries
layers so that no formal identity is mistaken for a uniform analytic theorem.

\section{Conclusions}

The Bell numbers admit two complementary complete descriptions.

First, on the principal Cauchy circle, they possess a rigorous fixed-order
expansion in $q=\W_0(n)/n$.  Every coefficient is a finite sum over partitions
of an even integer; its denominator and degree are controlled uniformly in the
order.  This resolves the usual ambiguity hidden behind phrases such as
``higher coefficients may be computed.''  The same formulas lead immediately
to factorial-free, logarithmic, and nested-log forms.

Second, the Hankel contour resolves the beyond-all-orders structure.  Every
Lambert branch $\W_k(n+1)$ gives a thimble sector with its own all-order
coefficient series.  The Bell contour fixes which sectors are active, and the
singulants order them exponentially.  For fixed $k\ne0$, the suppression
$\exp\{-2\pi^2k^2n/\log^2n\}$ explains why no finite algebraic expansion can
recover these oscillatory corrections.

Together, the principal partition calculus, the Lambert-polynomial branch
expansions, and the thimble sectors form a complete saddle-transseries
framework.  They also supply direct algorithms for extending the displayed
coefficients to any desired order without introducing unnamed recurrences or
undocumented shifts.

\appendix

\section{A sharper view of the contour remainder}

This appendix isolates the two inequalities used in
\cref{thm:principal}.

\begin{lemma}[Uniform local domination]\label{lem:local-domination}
There are absolute constants $\eta,c_0>0$ and $r_0$ such that, whenever
$r\ge r_0$, $n=r\ee^r$, and $|\theta|\le\eta/r$,
\begin{equation}\label{eq:local-domination-app}
 \RePart\Phi_r(\theta)\le-c_0n(1+r)\theta^2.
\end{equation}
\end{lemma}

\begin{proof}
The real part is an even analytic function of $\theta$, with value and first
derivative zero and second derivative $-A=-n(1+r)$.  For every fixed
$j\ge2$, the $2j$th derivative at the origin is bounded in magnitude by
$\ee^r\Touch_{2j}(r)\le\ee^r\Bell_{2j}r^{2j}$ for $r\ge1$.  If
$|\theta|\le\eta/r$, the sum of all normalized terms of order at least four is
bounded by $C\eta^2A\theta^2$ for a sufficiently small universal $\eta$; this
can also be seen directly from analyticity of
$\ee^{r\ee^{\ii\theta}}$ in the disk $|r\theta|<2\eta$.  Taking
$C\eta^2\le1/4$ gives the result, for example with $c_0=1/4$.
\end{proof}

\begin{lemma}[Outer-arc domination]\label{lem:outer-domination}
For the same fixed $\eta>0$, there is $c_1>0$ such that
\begin{equation}\label{eq:outer-domination-app}
 \sup_{\eta/r\le|\theta|\le\pi}\RePart\Phi_r(\theta)
 \le-c_1\frac{n}{r^2}
\end{equation}
for all sufficiently large $r$.
\end{lemma}

\begin{proof}
From \eqref{eq:global-real-bound} and
$1-\cos\theta\ge2\theta^2/\pi^2$,
\[
 \RePart\Phi_r(\theta)
 \le-\ee^r\left(1-
 \exp\left\{-\frac{2r\theta^2}{\pi^2}\right\}\right).
\]
On the stated arc the exponent is at least $2\eta^2/(\pi^2r)$.  Since
$1-\ee^{-x}\ge x/2$ for small positive $x$, the right-hand side is at most
$-c\ee^r/r=-cn/r^2$.
\end{proof}

\begin{proposition}[Two-scale remainder]\label{prop:two-scale-rem}
For every fixed $M$,
\begin{equation}\label{eq:two-scale-rem}
 R_M(n)=\bigO_M(q^M)
 +\bigO\!\left(\ee^{-c q^{-1/4}}\right)
 +\bigO\!\left(\ee^{-c n/r^2}\right).
\end{equation}
The two exponential terms are $\smallo(q^J)$ for every fixed $J$.
\end{proposition}

\begin{proof}
The first term is the integrated Taylor remainder on
$|u|\le q^{-1/8}$ (equivalent to $|u|\le t^{-1/4}$ above).  The second is the
Gaussian tail on the remainder of the local arc, and the third follows from
\cref{lem:outer-domination}.  Since
$q^{-1/4}\to\infty$ faster than $\log(1/q)$ and
$n/r^2\gg\log(1/q)$, both exponentials are beyond all powers of $q$.
\end{proof}

\section{Coefficient identities in expanded form}

For reference, the coefficient extractor may be written without auxiliary
polynomials as
\begin{align}
 C_m(r)
 &=\frac1{\sqrt{2\pi}}\int_{-\infty}^{\infty}
 \ee^{-u^2/2}
 [t^{2m}]\exp\!\left(
 \sum_{j\ge3}
 \frac{\Touch_j(r)(\ii u)^j}{j![r(1+r)]^{j/2}}t^{j-2}
 \right)\dd u.\label{eq:C-integral-app}
\end{align}
Likewise,
\begin{align}
 \widetilde C_m(w)
 &=\frac1{\sqrt{2\pi}}\int_{-\infty}^{\infty}
 \ee^{-u^2/2}
 [t^{2m}]\exp\!\left(
 \sum_{j\ge3}
 \frac{w^{j-1}+(-1)^j(j-1)!}
 {j!(1+w)^{j/2}}(\ii u)^jt^{j-2}
 \right)\dd u.\label{eq:Ctilde-integral-app}
\end{align}
At a fixed order, each coefficient extraction contains only finitely many
$j$; the infinite sums in \eqref{eq:C-integral-app} and
\eqref{eq:Ctilde-integral-app} are formal shorthand.

The reciprocal Stirling coefficients satisfy the useful recurrence
\begin{equation}\label{eq:reciprocal-stirling-recurrence}
 \widehat s_0=1,
 \qquad
 \widehat s_m=\frac1m\sum_{j=1}^{m}j a_j\widehat s_{m-j},
\end{equation}
where
\begin{equation}\label{eq:aj-reciprocal}
 a_j=
 \begin{cases}
 -\dfrac{B_{j+1}}{j(j+1)},&j\text{ odd},\\[0.6em]
 0,&j\text{ even}.
 \end{cases}
\end{equation}
The ordinary Stirling coefficients satisfy the same recurrence with the
opposite sign in \eqref{eq:aj-reciprocal}.

\section{Notation concordance}

\begin{center}
\begin{tabularx}{\textwidth}{@{}>{\raggedright\arraybackslash}p{0.18\textwidth}X@{}}
\toprule
symbol & meaning\\
\midrule
$\Bell_n$ & $n$th Bell number\\
$\Touch_n(x)$ & Touchard polynomial, so $\Bell_n=\Touch_n(1)$\\
$r$ & principal circle saddle $\W_0(n)$\\
$q$ & principal small parameter $r/n=\ee^{-r}$\\
$A$ & circle curvature $n(1+r)$\\
$C_m(r)$ & principal multiplicative saddle coefficient\\
$P_m(r)$ & numerator in $C_m=P_m/\{r^m(1+r)^{3m}\}$\\
$\Lambda_m(r)$ & logarithmic/cumulant coefficient\\
$s_m,\widehat s_m$ & ordinary and reciprocal Stirling coefficients\\
$N$ & shifted parameter $n+1$ in the Hankel representation\\
$w_k$ & Lambert saddle $\W_k(N)$\\
$\widetilde C_m(w)$ & arbitrary-thimble coefficient\\
$\cJ_k$ & oriented contribution of the $k$th thimble\\
$\cA_k$ & singulant comparing sector $k$ with sector $0$\\
$L,\ell$ & $\log n$ and $\log\log n$\\
\bottomrule
\end{tabularx}
\end{center}

\begin{thebibliography}{99}

\bibitem{Bell1934}
E.~T. Bell,
\newblock Exponential polynomials,
\newblock \emph{Annals of Mathematics} \textbf{35} (1934), 258--277.

\bibitem{Corless1996}
R.~M. Corless, G.~H. Gonnet, D.~E.~G. Hare, D.~J. Jeffrey, and D.~E. Knuth,
\newblock On the Lambert $W$ function,
\newblock \emph{Advances in Computational Mathematics} \textbf{5} (1996),
329--359.

\bibitem{deBruijn}
N.~G. de Bruijn,
\newblock \emph{Asymptotic Methods in Analysis},
\newblock 3rd ed., Dover, New York, 1981.

\bibitem{FlajoletSedgewick}
P.~Flajolet and R.~Sedgewick,
\newblock \emph{Analytic Combinatorics},
\newblock Cambridge University Press, Cambridge, 2009.

\bibitem{GrunwaldSerafin}
J.~Grunwald and G.~Serafin,
\newblock Explicit bounds for Bell numbers and their ratios,
\newblock \emph{Journal of Mathematical Analysis and Applications}
\textbf{549} (2025), Article 129527; arXiv:2408.14182.

\bibitem{Hayman1956}
W.~K. Hayman,
\newblock A generalisation of Stirling's formula,
\newblock \emph{Journal f\"ur die reine und angewandte Mathematik}
\textbf{196} (1956), 67--95.

\bibitem{Hwang2023}
H.-K. Hwang,
\newblock A curious identity arising from Stirling's formula and saddle-point
method on two different contours,
\newblock \emph{Electronic Journal of Combinatorics} \textbf{30} (2023),
Paper P4.3.

\bibitem{MoserWyman1955}
L.~Moser and M.~Wyman,
\newblock An asymptotic formula for the Bell numbers,
\newblock \emph{Transactions of the Royal Society of Canada, Section III}
\textbf{49} (1955), 49--54.

\bibitem{NISTW}
NIST Digital Library of Mathematical Functions,
\newblock \S4.13, Lambert $W$-function,
\newblock \url{https://dlmf.nist.gov/4.13}.

\bibitem{Paris2016}
R.~B. Paris,
\newblock The asymptotics of the Touchard polynomials,
\newblock \emph{Mathematica Aeterna} \textbf{6} (2016), 765--779;
arXiv:1606.07883.

\bibitem{Wojdylo2006}
J.~Wojdylo,
\newblock Computing the coefficients in Laplace's method,
\newblock \emph{SIAM Review} \textbf{48} (2006), 76--96.

\bibitem{ProveItCCC}
V.~Reshetnikov,
\newblock \emph{Combinatorial Coefficient Calculus}, ProveIt research draft,
\newblock \href{https://github.com/VladimirReshetnikov/ProveIt/tree/main/Analysis/FabiusFunction/docs/semi-formalized-research-frontiers/drafts/combinatorial-coefficient-calculus/Combinatorial_Coefficient_Calculus}{ProveIt repository path: combinatorial-coefficient-calculus}, accessed September 2026.

\bibitem{ProveItTransseries}
V.~Reshetnikov,
\newblock \emph{Series and Transseries}, ProveIt research drafts,
\newblock \href{https://github.com/VladimirReshetnikov/ProveIt/tree/main/Analysis/FabiusFunction/docs/semi-formalized-research-frontiers/drafts/series-and-transseries}{ProveIt repository path: series-and-transseries}, accessed September 2026.

\end{thebibliography}

\end{document}
